\documentclass[a4paper,oneside,12pt]{amsart}

%\usepackage[spanish, activeacute]{babel}
\usepackage[utf8]{inputenc}
%\usepackage[latin1]{inputenc}
\usepackage{latexsym}
\usepackage{multicol}
\usepackage{enumerate}
\usepackage{booktabs}
\usepackage[heightrounded]{geometry}
\geometry{top=2cm,bottom=2cm,left=2cm,right=2cm}
\usepackage{graphicx}

\DeclareMathOperator{\cond}{cond}
\DeclareMathOperator{\Id}{Id}
\DeclareMathOperator{\re}{Re}
%\DeclareMathOperator{\im}{Im}
\newcommand{\I}{\mathbb{I}}
\newcommand{\R}{\mathbb{R}}
\newcommand{\Q}{\mathbb{Q}}
\newcommand{\C}{\mathbb{C}}
\newcommand{\Z}{\mathbb{Z}}
\newcommand{\N}{\mathbb{N}}
\DeclareMathOperator{\im}{Im}

\newcommand{\nombremateria}
{Investigaci\'on Operativa}
\newcommand{\nombrecuatrimestre}
    {Segundo cuatrimestre de 2018}
\newcommand{\numeroparcial}
    { Segundo Parcial }
\newcommand{\fecha}
    {9 de Octubre de 2018}
\newcommand{\numerotema}
    {\bf 1}

%\oddsidemargin -1cm \evensidemargin -1cm \textwidth 17.5cm
%topmargin 1.2cm \textheight 25cm
%\setlength{\textheight}{25cm}

\def \ds{\displaystyle}
\newtheorem{quest}{Ejercicio}
\thispagestyle{empty}

\begin{document}

\hfill \hfill
\begin{tabular}[c]{|c|c|c|c|c|c|}
\hline 
1 & 2 & 3 & 4 & 5 & Calificaci\'on \\ \hline \hline
\hspace{1 cm} & \hspace{1 cm} & \hspace{1 cm} & \hspace{1 cm} & \hspace{1 cm} & \hspace{1.3 cm} \\
\hspace{1 cm} & \hspace{1  cm} & \hspace{1  cm} & \hspace{1  cm} & \hspace{1  cm} & \hspace{1.3 cm} \\ \hline

\end{tabular}

\vspace{-1cm}

\noindent Nombre y apellido:

\vspace{0.2cm}

\noindent Número de libreta:

\noindent \hrulefill 

\smallskip
\renewcommand{\baselinestretch}{1.1}

\begin{center}
	\large \textbf{\nombremateria} 
	
	\numeroparcial -- \fecha
\end{center}

\vspace{.6cm}
\renewcommand{\theenumi}{\textbf{\arabic{enumi}}}

%%% EMPIEZA EL EJERCICIO 1 %%%

\begin{quest}{\tiny .}



\begin{itemize}
    \item[(a)] Utilizar la primera fase del método de dos fases de SIMPLEX para  hallar una solución básica factible inicial de:
    \[\begin{array}{rrrrrrrr}
        \max & 2x_1 & + & 3x_2 & + & x_3  \\
        \text{s.a:} & x_1 & + & x_2 & + & x_3 & \leq & 40 \\
          & -2x_1 & - & x_2 & + & x_3 & \leq & -10 \\
          &   & - & x_2 & + & x_3 & \geq & 10 \\
         & \multicolumn{5}{r}{x_1,x_2,x_3} & \geq & 0
    \end{array}
    \]
    
    \item[(b)] Escribir el diccionario inicial de la Fase II. ¿Es la solución básica inicial encontrada en (a) la solución óptima del problema? Justificar. 
 
\end{itemize}

\end{quest}

%%% TERMINA EL EJERCICIO 1 %%%

%%% EMPIEZA EL EJERCICIO 2 %%%
\begin{quest}
El siguiente diccionario corresponde a una iteración de SIMPLEX para resolver un problema lineal de maximización:
\[
\begin{array}{rrrrrrrrr}
    x_6 & = & 4 & - & 7x_1 & + & 2x_3 & - & x_4  \\
    x_2 & = & 2+a & - & bx_1 &  & & - & 2x_4  \\
    x_5 & = & 5 & - & x_1 & + & cx_3 & - & x_4  \\ \hline
    z & = & 13 & + & dx_1 & + & ex_3 & - & 7x_4  \\
\end{array}
\]
Hallar condiciones sobre $a,b,c,d,e\in\mathbb{R}$ para que se cumpla que:
\begin{itemize}
    \item[(a)] el problema tiene al menos una solución factible, pero el problema dual asociado es infactible.
    \item[(b)] la próxima iteración es degenerada.
    \item[(c)] $x_6$ salga de la base en la próxima iteración
    \item[(d)] el valor óptimo del problema dual asociado sea 13.
\end{itemize}
\end{quest}
%%% TERMINA EL EJERCICIO 2 %%%

%%% EMPIEZA EL EJERCICIO 3 %%%


\begin{quest}
El siguiente modelo de programación lineal busca maximizar la ganancia obtenida al elaborar distintos tipos de productos, sujeto a la cantidad de materia prima y de horas hombre disponibles:
\[\begin{array}{rrrrrrrrc}
    \max & 6x_1 & + & 14x_2 & + & 13x_3 & & & (Ganancia)  \\
     & \frac{1}{2}x_1 & + & 2x_2 & + & x_3 & \leq & 24 & (Materia\; prima) \\
     & x_1 & + & 2x_2 & + & 4x_3 & \leq & 60 & (Horas \; hombre) \\
     & \multicolumn{5}{r}{x_1,x_2,x_3} & \geq & 0
\end{array}
\]
\begin{itemize}
    % \item[(a)] Plantear y resolver el problema dual.
    \item[(a)] Si se modifica en cantidades pequeñas la cantidad de materia prima disponible, ¿hasta cuánto conviene pagar por cada unidad extra de ella para aumentar la ganancia?
    \item[(b)] Resolver el problema para decidir cuánto elaborar de cada producto.
    \item[(c)] Hallar todos los $\alpha\in\mathbb{R}$ tales que la solución óptima obtenida en (b) siga siéndolo para el problema con función objetivo: $6x_1+(14+\alpha )x_2+13x_3$
    \item[(d)] ¿Cuánto pueden aumentarse las horas hombre de manera tal que la base óptima siga siéndolo?
\end{itemize}
\end{quest}

%%% TERMINA EL EJERCICIO 3 %%%

%%% EMPIEZA EL EJERCICIO 4 %%%


\begin{quest} 
Utilizar el algoritmo de Branch \& Bound para resolver el siguiente problema de programaci\'on lineal entera, detallando cada divisi\'on en subproblemas durante el algoritmo mediante el esquema de \'arbol visto en clase.   

\[\begin{array}{rrrrrrrrc}
\min & -4x_1 & + & 5x_2 &  &   \\
\text{s.a:}	 & 5x_1 & - & 2x_2 & \geq & -2  \\
			 & 2x_1 & - & 4x_2 & \leq & 5 \\
			 &	    &   & x_1 & \leq & 5 \\
			 &	    &   & x_2 & \leq & 4 \\
			 & \multicolumn{3}{r}{x_1,x_2} & \in & \mathbb{Z}_{\geq 0}
\end{array}
\]



\end{quest}
%%% TERMINA EL EJERCICIO 4 %%%

%%% EMPIEZA EL EJERCICIO 5 %%%


\begin{quest}
\textbf{(Sólo si debe recuperar el primer parcial)} Un avión de carga tiene tres compartimentos para transportar mercadería: uno en la parte delantera, uno en la parte central y otro en la parte trasera. Cada uno tiene la siguiente capacidad tanto de peso como de espacio:\\

\begin{tabular}{lcc} \toprule
    Compartimento & Capacidad de peso (Tn.) & Capacidad espacial ($m^3$)  \\ \midrule
    Delantero & 10 & 6800 \\
    Central & 16 & 8700 \\
    Trasero & 8 & 5300 \\ \bottomrule
\end{tabular} \\

Los siguientes tipos de mercancía están disponibles para el siguiente vuelo: \\

\begin{tabular}{lccc} \toprule
    Mercancía & Peso (Tn.) & Volumen ($m^3/Tn.$) & Ganancia ($\$/Tn.$) \\ \midrule
    $M_1$ & 18 & 480 & 310 \\
    $M_2$ & 15 & 650 & 380 \\
    $M_3$ & 23 & 580 & 350 \\
    $M_4$ & 12 & 390 & 285 \\ \bottomrule
\end{tabular} \\

Se puede transportar cualquier proporción de las mercancías, incluso podría no transportarse alguna de ellas. También se deben tener en cuenta los  siguientes requerimientos:
\begin{itemize}
    \item[(a)] El peso que carga el compartimento trasero debe ser proporcional al que carga el compartimento delantero.
    \item[(b)] La diferencia entre lo que se carga de $M_1$ y de $M_4$ en el compartimento central debe ser de al menos 3 Tn.
    \item[(c)] Al menos el $10\%$ de lo que se decida transportar de $M_2$ debe ser ubicado en el compartimento delantero.
    \item[(d)] En el compartimento trasero se puede ocupar una cantidad de espacio estrictamente menor que $1000m^3$ para $M_3$.
    \item[(e)] Si se transportan $6$ Tn. o más de $M_2$ entonces deben transportarse como mucho $16$ Tn. de $M_3$.
\end{itemize}

El objetivo es determinar cuánto de cada mercancía $M_1, M_2, M_3$ y $M_4$ debe transportarse y cómo distribuirla en los compartimentos del avión de manera tal que se maximice la ganancia. 
% La construcción de un barco requiere llevar a cabo las siguientes operaciones (la tabla también muestra cuántos días se requieren para cada operación):
% \begin{table}[htbp!]
%     \centering
%     \begin{tabular}{ccc} \toprule
%         Operaciones & Duración & Precedentes  \\ \midrule
%         A & 2 & ninguna\\ 
%         B & 4 & A \\ 
%         C & 2 & A \\ 
%         D & 5 & A \\ 
%         E & 3 & B o C \\ 
%         F & 3 & E \\ 
%         G & 2 & E \\ 
%         H & 7 & D, E y G \\ 
%         I & 4 & F o G \\ 
%     \end{tabular}
% \end{table}


% Algunas operaciones son alternativas a otras. En particular:
% \begin{enumerate}
%     \item sólo una entre B y C debe ser ejecutada
%     \item sólo una entre F y G debe ser realizada
% \end{enumerate}
% Además, 
% \begin{enumerate}
%     \setcounter{enumi}{2}
%     \item si se ejecutan las operaciones C y G, la duración de I incrementa 2 días
% \end{enumerate} 
% La tabla muestra también qué operaciones deben ser completadas antes del comienzo de una nueva operación. Por ejemplo, H puede comenzar sólo después de haber completado E, D y G (si G es ejecutada).\\
% Escribir un modelo de programación lineal que pueda ser utilizado para decidir qué operaciones ejecutar de manera tal que el tiempo de construcción del barco sea mínimo.
\end{quest}

%%% TERMINA EL EJERCICIO 5 %%%


%%% TERMINA EL PARCIAL %%%



\end{document}

